\documentclass[10pt,a4paper]{amsart}
\usepackage[margin=19mm]{geometry}
\usepackage{amsmath,amssymb,mathtools}
\usepackage[scr=rsfs,cal=euler]{mathalfa}
\usepackage{xfrac}
\usepackage[unicode,hidelinks,bookmarks=false]{hyperref}
\newcommand{\ie}{i.e.\ }
\newcommand{\cf}{cf.\ }
\newcommand{\resp}{respectively\ }
\newcommand{\Subsection}{\subsection}
\providecommand{\nobreakdash}{\nobreak\hbox{-}}
\setlength{\emergencystretch}{2em}
\multlinegap0pt
\usepackage{xcolor}
\usepackage{tikz}
\usepackage{amssymb}
\usepackage{dsfont}
\usepackage[shortlabels]{enumitem}
\newtheorem{lemma}{Lemma}
\newtheorem{theorem}{Theorem}
\usepackage{cleveref}
\crefname{lemma}{Lemma}{Lemmas}
\crefname{theorem}{Theorem}{Theorems}
\newcommand{\FR}[1]{\mathrm{FR}_{#1}}
\newcommand{\Flatten}[0]{\mathrm{Flatten}}
\newcommand{\Sym}{\mathfrak{S}}
\newcommand{\out}{\mathcal{O}}
\title{Counting $\ell$-Interval Fubini Rankings\\Through their Parking Outcomes}
\author{Bjorn Andreas Ager-Hart, Melissa Beerbower, Pamela E. Harris, Joakim Jakovleski, Matt McClinton}
\date{Source: arXiv:2603.00693v1 (CC BY 4.0)}
\begin{document}
\maketitle

\noindent\emph{Source and licence.} Counting $\ell$-Interval Fubini Rankings\\Through their Parking Outcomes. Version arXiv:2603.00693v1 (CC BY 4.0), distributed by arXiv under the Creative Commons Attribution 4.0 International licence, http://creativecommons.org/licenses/by/4.0/, as recorded in the arXiv licence metadata for this exact version and shown on the abstract page https://arxiv.org/abs/2603.00693v1. Full licence notice: \texttt{licenses/arxiv-2603.00693-CC-BY-4.0.txt}.\par\smallskip

\noindent\textit{Editorial scope.} Counted unit: the article's theorem outcome (source0.tex lines 699--702). The article's complete original proof of it is delivered with it (source0.tex lines 703--750, one \texttt{proof} environment of 48 lines). No mathematics is written, repaired or re-derived: the statement and the proof are the article's own lines, in the article's own notation, and no retained line is substituted or re-flowed.  Retained uncounted as the source-local context the proof needs: lines 626--629.  Nothing was omitted from any retained span, so the package carries no omission record.  No retained line contains a citation, so the wrapper carries no bibliography and no bibliography entry.  No retained line contains a figure environment, an image-inclusion command or drawing code, so no image is delivered and the package declares no asset dependency.  Editorial formatting only, in addition: an AMS \texttt{amsart} preamble that restates the article's own declarations and macros used by the retained lines, namely the environments \texttt{lemma}, \texttt{theorem} on separate counters, together with the standard packages those lines need.  The retained environments print in this document as Lemma 1, Theorem 1 in order of appearance, whereas the article prints the same items with the numbering of its own full document.  The complete native source of the article as distributed in the arXiv e-print for this exact version is bundled unchanged as \texttt{sources/195499/source0.tex}.
\begin{lemma}[Tied players park in sequential order]\label{lemma: competitor outcome for ties}
    Let $\alpha\in\FR{n}$ with $\psi(\alpha)=(B_1,B_2,\ldots,B_t)$.
    Then for each $i\in[t]$ the competitors in block $B_i$ appear in lexicographic order in the parking outcome $\pi = \out_{\FR{n}}(\alpha)$, in the spots numbered \[1+\sum_{j=1}^{i-1}|B_j|,~2+\sum_{j=1}^{i-1}|B_j|,~\ldots,~ \sum_{j=1}^{i}|B_j|. \]
\end{lemma}

\begin{theorem}\label{thm:ordered set partition is the outcome}
    If $\alpha\in\FR{n}$, then 
    $\out_{\FR{n}}(\alpha)=\Flatten(\psi(\alpha))$.
\end{theorem}

\begin{proof}
    Let $\alpha\in\FR{n}$ and let $\psi(\alpha)=(B_1,B_2,\ldots,B_t)$ be its corresponding ordered set partition, where we order the entries in each block in increasing order. 
    We also recall that each block of competitors corresponds to a rank in $\alpha$, and the blocks are sorted by increasing rank. Namely, if competitor $b_i$ is in block $B_i$, competitor $b_j$ is in block $B_j$, and $i < j$, then $a_{b_i} < a_{b_j}$ i.e., competitor $b_i$ ranks \textit{before} competitor $b_j$. 
    For each $i\in[t]$, label the elements in each block as $B_i=\{b_{(i,1)},b_{(i,2)},\ldots,b_{(i,|B_i|)}\}$, where we assume that $b_{(i,j)}<b_{(i,j+1)}$ for all $j\in[|B_i|-1]$. 
    Create the word $w=w_1w_2\cdots w_n=\Flatten(\psi(\alpha))$.
    
    
    Let
    $\out_{\FR{n}}(\alpha)=\pi=\pi_1\pi_2\cdots\pi_n\in\Sym_n$. 
    For each $i\in [n]$, the value $\pi_i$ is the label of the car parked in spot~$i$.
    We want to show that 
    $\pi_i=w_i$, for all $i\in[n]$.
    
    
    We proceed by induction on the index $i\in[n]$. 
    To start, recall that by \Cref{lemma: competitor outcome for ties} 
    the cars in block 
    $B_1=\{b_{(1,1)},b_{(1,2)},\ldots,b_{(1,|B_1|)}\}$ park in spots $1,2,\ldots,|B_1|$.
    Thus, $\pi_i=b_{(1,i)}$ for each $i\in[|B_1|]$. 
    Moreover, by construction of $w$, we have $w_i=b_{(1,i)}$ for each $i\in[|B_1|]$. By transitivity, $\pi_i=w_i$ for all $i\in[|B_1|]$.
    This establishes the base case of the induction on the first $|B_1|$ cars.
    
    Assume for induction that $\pi_i=w_i$ for all $i< n$.
    Now consider $\pi_{i+1}$ and $w_{i+1}$.
    Given $w_i$, we know there exists $j\in[t]$ and $s\in[|B_j|]$ such that 
    $w_{i}=b_{(j,s)}$, i.e., $w_i$ is the $s$-th entry in block $B_j$. 
    Then there are two possibilities for what $w_{i+1}$ could be:
    \begin{itemize}
        \item $w_{i+1}=b_{(j,s+1)}$ if $w_i$ was \textit{not} the last entry of a block, and 
        \item $w_{i+1}=b_{(j+1,1)}$ if $w_i$ was the last entry of a block.
    \end{itemize}
    In the case when $w_{i+1}=b_{(j,s+1)}$, we know by \Cref{lemma: competitor outcome for ties} that car $b_{(j,s)}$ parks in spot 
    $y=s+\sum_{x=1}^{j-1}|B_x|$,
    and so car $b_{(j,s+1)}$ parks in spot 
    $y+1=s+1+\sum_{x=1}^{j-1}|B_x|$.
    Hence $\pi_y=b_{(j,s)}$ and $\pi_{y+1}=b_{(j,s+1)}$.
    
    Now recall that, by assumption, we have $w_i=b_{(j,s)}$, which means that $i=s+\sum_{x=1}^{j-1}|B_x|$. Thus, $i=y$. 
    So by transitivity,  this establishes that 
    $\pi_{i+1}=w_{i+1}$, as desired.
    
    In the case when $w_{i+1}=b_{(j+1,1)}$, we know that $w_{i+1}$ is the first entry in block $B_{j+1}$ and, by \Cref{lemma: competitor outcome for ties}, car $b_{(j+1,1)}$ parks in spot 
    $z=1+\sum_{x=1}^j|B_x|$.
    Thus $\pi_{z}=b_{(j+1,1)}$.
    Now recall that by assumption $w_{i+1}=b_{(j+1,1)}$, which means that $i+1=1+\sum_{x=1}^{j}|B_x|$. Thus, $i+1=z$.  
    So by transitivity, this yields 
    $\pi_{i+1}=w_{i+1}$, which completes the proof.
\end{proof}

\end{document}
